% TODOS
\newcommand{\kb}[1]{\todo[color=yellow,size=\tiny,fancyline,author=Kevin]{#1}\xspace}
\newcommand{\kbin}[1]{\todo[inline,color=yellow,author=Kevin]{#1}\xspace}
\newcommand{\atkb}{\colorbox{yellow}{\textbf{@Kevin}}}



\newcommand{\Nats}{\mathbb{N}}
\newcommand{\Reals}{\mathbb{R}}
\newcommand{\PosRealsInf}{\Reals^{\infty}}

% Macros for language keywords to be used in math mode
\newcommand{\letkw}{\textnormal{\ttfamily\bfseries let}}
\newcommand{\inkw}{\textnormal{\ttfamily\bfseries in}}
\newcommand{\ifkw}{\textnormal{\ttfamily\bfseries if}}
\newcommand{\thenkw}{\textnormal{\ttfamily\bfseries then}}
\newcommand{\elsekw}{\textnormal{\ttfamily\bfseries else}}
\newcommand{\flip}{\textnormal{\ttfamily\bfseries flip}} % Shortcut for uniform(0,1) < 0.5
\newcommand{\uniform}{\textnormal{\ttfamily\bfseries uniform}}
\newcommand{\discrete}{\textnormal{\ttfamily\bfseries discrete}}
\newcommand{\gaussian}{\textnormal{\ttfamily\bfseries gaussian}}
\newcommand{\exponential}{\textnormal{\ttfamily\bfseries exponential}}
\newcommand{\betafn}{\textnormal{\ttfamily\bfseries beta}} % Use beta for the keyword
\newcommand{\fstkw}{\textnormal{\ttfamily\bfseries fst}}
\newcommand{\sndkw}{\textnormal{\ttfamily\bfseries snd}}
\newcommand{\funkw}{\textnormal{\ttfamily\bfseries fun}}
\newcommand{\observekw}{\textnormal{\ttfamily\bfseries observe}}
\newcommand{\matchkw}{\textnormal{\ttfamily\bfseries match}}
\newcommand{\withkw}{\textnormal{\ttfamily\bfseries with}}
\newcommand{\refkw}{\textnormal{\ttfamily\bfseries ref}}
\newcommand{\fixkw}{\textnormal{\ttfamily\bfseries fix}}
\newcommand{\matchcase}{\boldsymbol{\mid}} % The | in the match e with ...
\newcommand{\diverge}{\textnormal{\ttfamily\bfseries diverge}}
\newcommand{\logand}{\mathrel{\&\&}}                    % logical and
\newcommand{\logor}{\mathrel{||}}                       % logical or
\newcommand{\lognot}{\text{not}}                        % logical not

% Types
\newcommand{\bool}{\textnormal{\ttfamily\bfseries bool}}
\newcommand{\intty}{\textnormal{\ttfamily\bfseries int}}
\newcommand{\float}{\textnormal{\ttfamily\bfseries float}}
\newcommand{\floattype}[2]{\mathnormal{\float}[#1; #2]}
\newcommand{\fin}[1]{\textnormal{\ttfamily\bfseries fin}(#1)}
\newcommand{\listty}[1]{\textnormal{\ttfamily\bfseries list}(#1)}
\newcommand{\unit}{\textnormal{\ttfamily\bfseries unit}}

% Macros for finite type operations
\newcommand{\finconst}[2]{#1_{\##2}}                    % k_{#n}
\newcommand{\finlt}[1]{<_{\#{#1}}}                      % <#n with subscript
\newcommand{\finleq}[1]{\leq_{\#{#1}}}                  % <=#n with subscript
\newcommand{\fingt}[1]{>_{\#{#1}}}                      % >#n with subscript
\newcommand{\fingeq}[1]{\geq_{\#{#1}}}                  % >=#n with subscript
\newcommand{\fineq}[1]{==_{\#{#1}}}                     % ==#n with subscript

% Macros for float type notation with proper braces (to be used in lstlisting)
\definecolor{typecolor}{HTML}{4169E1} % RoyalBlue
\newcommand{\codetype}[1]{\textcolor{typecolor}{\ttfamily\bfseries#1}}
\newcommand{\ftnv}[1]{\codetype{float[#1]}}  % Cuts only, no values
\newcommand{\ftb}[2]{\codetype{float[#1; #2]}}  % Both bags have values
\newcommand{\ft}[1]{\codetype{float[#1; T]}}        % Cuts only, values are T
\newcommand{\ftv}[1]{\codetype{float[T; #1]}}       % Cuts are T, values specified
\newcommand{\ftt}{\codetype{float[T; T]}}               % Both bags are T

\newcommand{\subst}[2]{[#1/#2]}
% Meta variables
\newcommand{\expr}{e}
\newcommand{\val}{v}
\newcommand{\var}{x}

\newcommand{\true}{\textnormal{true}}
\newcommand{\false}{\textnormal{false}}

% Type system
\newcommand{\lattice}[1]{\mathcal{L}(#1)}
\newcommand{\latticedom}{\mathcal{D}}
\newcommand{\latticeleq}{\sqsubseteq}
\newcommand{\latticesup}{\sqcup}
\newcommand{\latticeelem}{d}

\newcommand{\cutsets}{\mathcal{B}}
\newcommand{\clt}[1]{<\!\!{#1}}
\newcommand{\cleq}[1]{\leq\!\!{#1}}


\newcommand{\valsets}{\mathcal{V}}
\newcommand{\distinguishes}[2]{\mathsf{recovers}(#1, #2)}
\newcommand{\answerslt}[3]{\mathsf{answers}_<(#1, #2, #3)}
\newcommand{\answersleq}[3]{\mathsf{answers}_{\leq}(#1, #2, #3)}
\newcommand{\subtype}[2]{#1 <: #2}
\newcommand{\has}[2]{\mathsf{has}(#1, #2)}

\newcommand{\intervals}[1]{\text{Intervals}(#1)}


% Semantics
\newcommand{\dist}[1]{\mathcal{D}\left( #1 \right)}
\newcommand{\probmeas}{\mu}
\newcommand{\discup}{\mathbin{\dot{\cup}{}}}
\newcommand{\exprs}{\textnormal{Expr}}
\newcommand{\extexprs}{{\overline{\textnormal{Expr}}}}
\newcommand{\distextexprs}{\mathcal{D} \big( \extexprs \big)}
\newcommand{\distexprs}{\dist{\exprs}}
\newcommand{\obsfail}{\textnormal{\Lightning}}
\newcommand{\error}{\textnormal{error}}
\newcommand{\sem}[1]{\llbracket #1  \rrbracket}
\newcommand{\bigsemstep}[2]{\langle #1  \rangle_{#2}}
\newcommand{\bigsemstepterm}[2]{\langle #1  \rangle_{#2}^{\downarrow}}
\newcommand{\bigsem}[1]{\langle #1  \rangle}
\newcommand{\measset}{A}
\newcommand{\nonvals}{N}
\newcommand{\vals}{Vals}

\newcommand{\monunit}{\textnormal{return}}
\newcommand{\monbind}{\ensuremath{\ensuremath{\gg\!\!=}}}
\newcommand{\dirac}[1]{\delta_{#1}}



% Macros for language names, use small caps
\newcommand{\Slice}{\textnormal{\scshape Slice}\xspace}
\newcommand{\Storm}{\textnormal{\scshape Storm}\xspace}
\newcommand{\Dice}{\textnormal{\scshape Dice}\xspace}
\newcommand{\Roulette}{\textnormal{\scshape Roulette}\xspace}
\newcommand{\SPPL}{\textnormal{\scshape Sppl}\xspace}
\newcommand{\PSI}{\textnormal{\scshape Psi}\xspace}

\newcommand{\R}{\mathbb{R}}

%discretization
\newcommand{\discretize}[1]{\mathcal{D}\sem{#1}}

% DiscreteCase for the semantics
\newcommand{\dc}{\ensuremath{\texttt{DC}_{B,B_1,B_2,V_1,V_2}}}

% coupling relation
\newcommand{\relation}{\mathcal{R}}

% Macros for exact inference outputs
\newcommand{\result}[2]{\{\textit{true}: #1, \textit{false}: #2\}}

% Macros for expression skeletons
\newcommand{\skel}{Skel}
\newcommand{\holes}{holes}
\newcommand{\decompose}{\mathsf{dec}}

% Colors 
\colorlet{green}{green!70!black}
\colorlet{gray}{gray!70!black}
\colorlet{blue}{blue!50!black}
\colorlet{red}{red!80!black}
\colorlet{purple}{purple!80!black}
\colorlet{orange}{orange}
\colorlet{teal}{teal!80!black}
\colorlet{olive}{olive!70!black}
\colorlet{brown}{brown!90!black}
\colorlet{violet}{violet!70!black}

% Enumerate item
\newcommand*\circled[1]{\tikz[baseline=(char.base)]{
            \node[shape=circle,draw,inner sep=1.5pt] (char) {#1};}}
